\begin{lemma}
Let $g\in\mathsf{IncSeq}$ and let $k\in\mathbb N$ satisfy
$k<g(k)$ and $g(j)=j$ for every $j<k$.  For every
$f\in\mathsf{IncSeq}$ and every $l\in\mathbb N$, define
$\sigma_f(l)=\mathsf{BlockSigma}(g,k,f,l)$ as in
\noderef{BlockSigma}.  The block maps satisfy
\[
\sigma_{f\circ\mathsf{SuccSeq}}(l)=\sigma_f(g(l)).
\]
Here $f\circ\mathsf{SuccSeq}$ is the right composition represented by
$\mathsf{RightComp}(\mathsf{SuccSeq},f)$, and equivalently
$\sigma_{f\circ\mathsf{SuccSeq}}=\sigma_f\circ g$.
\end{lemma}
\begin{proof}
Put $G(n)=g^{[n]}(k)$, so $G(0)=k$.  The two cases below exhaust all
$l$ by \noderef{OrbitBlockCoverage}.  If $l<G(0)$, then $l<k$ and the
least-moved-point hypothesis gives $g(l)=l$.  By the piece of the block
definition in \noderef{BlockSigma} below $G(0)$, both sides are therefore
$l$.

Otherwise, let $n$ be the unique index such that
$G(n)\leq l<G(n+1)$.  Since $g$ is strictly increasing,
\[
G(n+1)=g(G(n))\leq g(l)<g(G(n+1))=G(n+2).
\]
Thus $g(l)$ is in the $(n+1)$st block.  Using the block formula from
\noderef{BlockSigma} and the composition convention from
\noderef{RightComp} and \noderef{SuccSeq}, the left-hand side is
\[
g^{[f(n+1)-n]}(l),
\]
while the right-hand side is
\[
g^{[f(n+1)-(n+1)]}(g(l)).
\]
The embedding $f$ satisfies $n+1\leq f(n+1)$, so with
$a=f(n+1)-(n+1)$ we have $a+1=f(n+1)-n$.  Hence the elementary iterate
identity $g^{[a]}(g(l))=g^{[a+1]}(l)$ makes these two values equal.  This
proves the claimed identity in both cases.
\end{proof}
